\section{This Paper, Generated by this Process}
\label{sec:self}

The fifth integrated practice (\S\ref{sec:practice-recursive})
requires that a paper which prescribes a discipline have been written
under that discipline. This section is the explicit attestation: the
manuscript you are reading was produced by exactly the
F(AI)\textsuperscript{2}R pipeline it describes. The attestation is
verifiable; every claim below points to an artefact in the public
repository.

\subsection{What ran}

The repository\footnote{\url{https://github.com/noheton/f-ai-r}} is the
worked example. The ten-stage pipeline of \S\ref{sec:pipeline} is
materialised as one markdown file per role under \texttt{agents/},
each file structured per the prompt taxonomy in
Appendix~\ref{app:prompts}. The Orchestrator routes work; the
Scientific Writer drafts prose; the Source Analyzer upgrades retrieved
sources; the Layout Scrutinizer, Readability \& Novelty Scrutinizer,
and Aligner file defect registries; the Provenance Curator alone
writes \texttt{doc/provenance.ttl}; the Site Agent renders the public
site. The harness was Anthropic Claude
(\texttt{claude-opus-4-7}) operating under the rules in
\texttt{CLAUDE.md}.

\subsection{What is committed}

\begin{itemize}
    \item \textbf{The manuscript.} \texttt{paper/main.tex} plus the
        per-section files under \texttt{paper/sections/}, one
        \texttt{\textbackslash section\{\}} per file
        (\S\ref{sec:practice-mirror}).
    \item \textbf{The agents.} Ten prompt files under
        \texttt{agents/}, each carrying the binding
        primary-artifact-consistency block of \S\ref{sec:mirror}.
    \item \textbf{The provenance graph.}
        \texttt{doc/provenance.ttl}, valid Turtle, parsed by
        \texttt{rdflib} as part of the CI gate. Every named claim in
        this manuscript has a corresponding
        \texttt{fair2r:Claim} IRI in the graph.
    \item \textbf{The logbook.} \texttt{doc/logbook.md}, append-only,
        ISO-8601-dated, one entry per working session.
    \item \textbf{The build pipeline.}
        \texttt{.github/workflows/build-paper.yml} compiles both
        manuscripts on every push;
        \texttt{.github/workflows/pages.yml} validates the Turtle,
        builds the static site, and (on \texttt{main}) deploys to
        GitHub Pages.
\end{itemize}

\subsection{How a third party can replay it}

\begin{enumerate}
    \item \texttt{git clone https://github.com/noheton/f-ai-r}
    \item Install a TeX Live distribution
        (\texttt{latexmk}, \texttt{biber}, \texttt{lmodern},
        \texttt{biblatex}) and the dependencies in
        \texttt{scripts/requirements.txt}.
    \item \texttt{make -C paper all} produces \texttt{paper/main.pdf}
        and \texttt{paper/main-condensed.pdf}.
    \item \texttt{python scripts/build\_provenance\_site.py} produces
        \texttt{\_site/}, including the rendered provenance tables and
        the agent prompts.
    \item Compare against the published artefacts; any difference is a
        defect.
\end{enumerate}

\subsection{Honest limits of the attestation}

The attestation is artefactual, not behavioural. It demonstrates that
the eight practices and the ten-stage pipeline produced an artefact
set that looks the way the prescription says it should look. It does
not, by itself, demonstrate that the human author resisted every
shortcut the pattern was meant to prevent --- that depends on review.
The repository is what it is; the prose is what it is; the graph
either matches or it does not.

\paragraph{What this section is.}
A signed claim of recursive practice. The signing happens not in this
prose but in the commits that produced it: every paragraph here is a
\texttt{fair2r:Claim} in \texttt{doc/provenance.ttl}, attributed to
the human author, generated by an authoring activity, supported (where
appropriate) by a vendored source. \emph{If the graph and the prose
disagree, the graph wins, and you have just found a defect we owe you
a fix for.}
